\chapter{Two Tomes at a Time}

\begin{center}\itshape
The art of doing mathematics consists in finding that special case which contains all the germs of generality.\\
\upshape --- David Hilbert
\end{center}

In 1977 a letter reached Martin Gardner, who wrote the ``Mathematical Games'' column of \emph{Scientific American}. Its author,
George Sicherman of Buffalo, had found two strange dice. The first carries the numbers $1$, $2$, $2$, $3$, $3$, $4$ on its faces, the second
$1$, $3$, $4$, $5$, $6$, $8$. Throw them together and add: the sum is $2$ with probability $\frac1{36}$, $7$ with probability $\frac6{36}$,
$12$ with probability $\frac1{36}$, and every other sum has exactly the same chance as with an ordinary pair
(Figure~\ref{fig:dice}). A gambler could not tell them apart by watching the totals. Gardner published them in February 1978,
and asked the natural question: are there any others?

The question is about chance, and the answer is about polynomials. Write the faces of a die as the exponents of a polynomial,
$x+x^2+\dots+x^6$ for an ordinary die; then throwing two dice and adding multiplies the polynomials, and the question becomes:
in how many ways can $(x+\dots+x^6)^2$ be split into two factors that look like dice? That is Problem~\ref{pb:sicherman}, and
it needs Tome~3 to set it up and Tome~0 to finish it. This chapter is made of such pairs. Each problem crosses from one volume
to another, and most are short; they are the warm-up for the rest of the book.

\begin{think}
Before reading on, try to find a third pair of dice with the same sums, or convince yourself that there is none. What must the
two smallest faces be? And could two \emph{loaded} dice, with faces $1$ to $6$, make all eleven sums from $2$ to $12$ equally
likely?
\end{think}

\begin{figure}[ht]
\centering
\fig{ch01-dice}
\caption{The sums of two ordinary dice (left) and of Sicherman's dice $1$, $2$, $2$, $3$, $3$, $4$ and $1$, $3$, $4$, $5$, $6$, $8$ (right). The colour
marks the sum: each colour covers the same number of cells in both tables.}
\label{fig:dice}
\end{figure}

\section{Dice, polynomials and primes}

A polynomial is a list of numbers that has learned to multiply. The first problems use this in both directions: chance and
arithmetic questions become questions about factoring polynomials, and polynomials over $\mathbb F_p$ turn out to be nothing
more than functions.

\begin{exercise}[$\star\star$]\label{pb:sicherman}
Find all pairs of dice with six faces each, every face carrying a positive integer (repetitions allowed), such that the sum
of the two dice, thrown independently, has the same law as the sum of two ordinary dice.
\end{exercise}

\begin{exercise}[$\star\star$]\label{pb:loaded}
Two dice have faces $1,2,\dots,n$, with $n\ge2$, but they are loaded: the first shows $i$ with probability $p_i$, the
second $j$ with probability $q_j$, independently. Show that the sum cannot be uniformly distributed on $\{2,3,\dots,2n\}$.
Give two proofs for $n=6$: one with the roots of a polynomial, one with an inequality.
\end{exercise}

\begin{exercise}[$\star\star$]\label{pb:primepoly}
Let $P\in\Z[X]$ be non-constant.
\begin{enumerate}[label=(\alph*)]
\item Show that there are infinitely many integers $n\ge0$ for which $|P(n)|$ is not a prime number.
\item Show that infinitely many primes divide at least one of the numbers $P(0),P(1),P(2),\dots$
\end{enumerate}
\end{exercise}

\begin{exercise}[$\star\star$]\label{pb:fpfunctions}
Let $p$ be a prime.
\begin{enumerate}[label=(\alph*)]
\item Show that every map $f:\mathbb F_p\to\mathbb F_p$ is the map $x\mapsto P(x)$ of exactly one polynomial $P\in\mathbb F_p[X]$
of degree at most $p-1$.
\item Show that the coefficient of $X^{p-1}$ in this polynomial is $-\sum_{a\in\mathbb F_p}f(a)$.
\item Suppose $p$ odd and $f$ a bijection. Show that $P$ has degree at most $p-2$, and that for each $t$ with $1\le t\le p-2$
the polynomial of the map $x\mapsto f(x)^t$ has degree at most $p-2$ too. Check this for $f(x)=x^3$ on $\mathbb F_5$.
\end{enumerate}
\end{exercise}

\section{Old friends in new fields}

Fibonacci numbers, binomial coefficients and Catalan numbers are the first sequences one meets. Here they are taken to places
where they were not born: finite fields, the $p$-adic numbers, the binary digits of their indices.

\begin{exercise}[$\star\star$]\label{pb:fibmodp}
Let $p\ne2,5$ be a prime and $\left(\frac5p\right)$ the Legendre symbol. Let $(F_n)$ be the Fibonacci numbers.
\begin{enumerate}[label=(\alph*)]
\item Show that $p$ divides $F_{p-(5/p)}$: so $p\mid F_{p-1}$ when $p\equiv\pm1\pmod5$, and $p\mid F_{p+1}$ when
$p\equiv\pm2\pmod5$.
\item Show that the period of $(F_n\bmod p)$ divides $p-1$ in the first case and $2(p+1)$ in the second.
\item Find the periods for $p=7$ and $p=11$.
\end{enumerate}
\end{exercise}

\begin{exercise}[$\star\star$]\label{pb:whichroot}
For $n\ge0$, $\binom{1/2}n=\frac1{n!}\cdot\frac12\bigl(\frac12-1\bigr)\cdots\bigl(\frac12-n+1\bigr)$.
\begin{enumerate}[label=(\alph*)]
\item Show that $\binom{1/2}n=(-1)^{n-1}\frac{C_{n-1}}{2^{2n-1}}$ for $n\ge1$, where $C_m=\frac1{m+1}\binom{2m}m$ is the $m$-th
Catalan number.
\item Show that $\sum_{n\ge0}\binom{1/2}n7^n$ converges in $\Q_7$ to a square root of $8$, and say which one.
\item Show that $\sum_{n\ge0}\binom{1/2}n8^n$ converges in $\Q_2$, and that its sum is $-3$.
\end{enumerate}
\end{exercise}

\begin{exercise}[$\star$]\label{pb:pascal2}
Let $s(n)$ be the number of $1$s in the binary expansion of $n$.
\begin{enumerate}[label=(\alph*)]
\item Show that $\binom{2^k+m}j\equiv\binom mj+\binom m{j-2^k}\pmod2$ for $0\le m<2^k$ (with $\binom mi=0$ for $i<0$),
and explain the pattern of Figure~\ref{fig:pascal}.
\item Deduce that row $n$ of Pascal's triangle contains exactly $2^{s(n)}$ odd numbers.
\item Show that the rows $0,1,\dots,2^k-1$ contain $3^k$ odd numbers in all, and that the proportion of odd entries in them
tends to $0$.
\end{enumerate}
\end{exercise}

\begin{figure}[ht]
\centering
\fig{ch01-pascal}
\caption{The odd entries of the first $32$ rows of Pascal's triangle. Rows $16$ to $31$ repeat rows $0$ to $15$ twice, side
by side (Problem~\ref{pb:pascal2}).}
\label{fig:pascal}
\end{figure}

\begin{exercise}[$\star\star$]\label{pb:millin}
\begin{enumerate}[label=(\alph*)]
\item Show that $F_{2m}=F_m(F_{m-1}+F_{m+1})$ and $F_{2m-1}=F_m^2+F_{m-1}^2$ for $m\ge1$.
\item Show that $\sum_{n=0}^N\frac1{F_{2^n}}=3-\frac{F_{2^N-1}}{F_{2^N}}$ for $N\ge1$, and deduce that
\[
\sum_{n\ge0}\frac1{F_{2^n}}=1+1+\frac13+\frac1{21}+\frac1{987}+\cdots=\frac{7-\sqrt5}2.
\]
\end{enumerate}
\end{exercise}

\begin{exercise}[$\star\star$]\label{pb:sylvester}
\begin{enumerate}[label=(\alph*)]
\item (Fibonacci, 1202.) Let $0<\frac ab<1$ be rational. Subtract the largest unit fraction $\frac1n\le\frac ab$, and repeat
with the remainder. Show that the process stops, and writes $\frac ab$ as a sum of distinct unit fractions. Apply it to
$\frac4{13}$.
\item (Sylvester, 1880.) Let $s_0=2$ and $s_{k+1}=s_k^2-s_k+1$, so $2$, $3$, $7$, $43$, $1807$, \dots Show that
$\sum_{k=0}^N\frac1{s_k}=1-\frac1{s_{N+1}-1}$, and that the $s_k$ are pairwise coprime.
\item Show that there is a constant $E>1$ such that $s_k=\bigl\lfloor E^{2^{k+1}}+\frac12\bigr\rfloor$ for every $k\ge0$.
\end{enumerate}
\end{exercise}

\section{Arithmetic functions in disguise}

Euler's $\varphi$, Möbius's $\mu$ and von Mangoldt's $\Lambda$ were born in Tome~1 as tools for counting. They reappear as
determinants, as coefficients of power series, and as the exact value of an infinite sum.

\begin{exercise}[$\star\star$]\label{pb:smith}
(Smith, 1875.)
\begin{enumerate}[label=(\alph*)]
\item Show that $\det\bigl(\gcd(i,j)\bigr)_{1\le i,j\le n}=\varphi(1)\varphi(2)\cdots\varphi(n)$.
\item More generally, let $f$ be any function on the positive integers and $g=f*\mathbf1$, that is $g(m)=\sum_{d\mid m}f(d)$.
Show that $\det\bigl(g(\gcd(i,j))\bigr)_{1\le i,j\le n}=f(1)\cdots f(n)$, and compute $\det\bigl(\gcd(i,j)^2\bigr)_{1\le i,j\le n}$.
\end{enumerate}
\end{exercise}

\begin{exercise}[$\star\star$]\label{pb:lambert}
(Lambert series.) Let $(a_n)_{n\ge1}$ be complex numbers with $|a_n|\le Cn^K$ for some constants $C,K$.
\begin{enumerate}[label=(\alph*)]
\item Show that for $|x|<1$, $\displaystyle\sum_{n\ge1}a_n\frac{x^n}{1-x^n}=\sum_{m\ge1}\Bigl(\sum_{d\mid m}a_d\Bigr)x^m$.
\item Deduce the sums of $\sum_n\varphi(n)\frac{x^n}{1-x^n}$, $\sum_n\mu(n)\frac{x^n}{1-x^n}$ and $\sum_n\lambda(n)\frac{x^n}{1-x^n}$,
where $\lambda(n)=(-1)^{\Omega(n)}$ is Liouville's function ($\Omega(n)$ counts the prime factors of $n$ with multiplicity).
\item Show that $\displaystyle\sum_{n\ge1}\frac{\varphi(n)}{2^n-1}=2$ and $\displaystyle\sum_{n\ge1}\frac{\mu(n)}{2^n-1}=\frac12$.
\end{enumerate}
\end{exercise}

\begin{exercise}[$\star\star$]\label{pb:coprime52}
Show that
\[
\sum_{\substack{a,b\ge1\\\gcd(a,b)=1}}\frac1{a^2b^2}=\frac52.
\]
More generally, compute $\sum_{\gcd(a,b)=1}(ab)^{-s}$ for $s>1$, and interpret the result as a Dirichlet series
$\sum_nc_nn^{-s}$ with an explicit $c_n$.
\end{exercise}

\section{Chance meets structure}

Probability is full of hidden algebra: a random walk on a polygon hides the roots of unity, a moment of a random permutation
hides an orbit count, the carries of an addition hide a Markov chain. The problems of this section uncover them.

\begin{exercise}[$\star$]\label{pb:cyclewalk}
A token moves on the vertices $0,1,\dots,k-1$ of a regular polygon ($k\ge3$); at each step it moves to one of the two
neighbouring vertices with probability $\frac12$ each, independently of the past. Let $q_n$ be the probability that it is at
$0$ after $n$ steps, having started at $0$.
\begin{enumerate}[label=(\alph*)]
\item Show that $q_n=\frac1k\sum_{j=0}^{k-1}\cos^n\frac{2\pi j}k$.
\item Give $q_n$ in closed form for $k=3$, $4$ and $6$.
\item Find the limit of $q_n$ as $n\to\infty$, according to the parities of $k$ and $n$.
\end{enumerate}
\end{exercise}

\begin{exercise}[$\star\star$]\label{pb:fixmoments}
Let $F$ be the number of fixed points of a uniformly random permutation $\sigma$ of $\{1,\dots,n\}$, and let $\mathfrak S_n$
act on $\{1,\dots,n\}^k$ by $\sigma\cdot(x_1,\dots,x_k)=(\sigma(x_1),\dots,\sigma(x_k))$.
\begin{enumerate}[label=(\alph*)]
\item Show that $\mathbb E[F^k]$ is the number of orbits of this action.
\item Deduce that $\mathbb E[F^k]=\sum_{j=0}^n\stirlingii kj$, and that $\mathbb E[F^k]$ is the Bell number $B_k$ when $k\le n$.
\item Deduce that the first $n$ moments of $F$ are exactly those of a Poisson variable of parameter $1$.
\end{enumerate}
\end{exercise}

\begin{exercise}[$\star\star$]\label{pb:matrices2}
Let $p$ be a prime, and let $M$ be a uniformly random matrix of $M_2(\mathbb F_p)$.
\begin{enumerate}[label=(\alph*)]
\item Count the matrices that are scalar; that have two distinct eigenvalues in $\mathbb F_p$; whose characteristic polynomial
is irreducible over $\mathbb F_p$; that are not scalar but have a double eigenvalue.
\item Deduce that, as $p\to\infty$, $M$ is diagonalisable over $\mathbb F_p$ with probability tending to $\frac12$, and has an
irreducible characteristic polynomial with probability tending to $\frac12$.
\end{enumerate}
\end{exercise}

\begin{exercise}[$\star\star$]\label{pb:carries}
Two numbers with $n$ digits in base $b\ge2$ are chosen at random, all $2n$ digits being independent and uniform on
$\{0,\dots,b-1\}$, and added by hand. Let $C_i$ be the carry out of position $i$ (position $1$ is the units), and $C_0=0$.
\begin{enumerate}[label=(\alph*)]
\item Show that $(C_i)_{i\ge0}$ is a Markov chain on $\{0,1\}$, and find its transition matrix and its eigenvalues.
\item Show that $P(C_i=1)=\frac12(1-b^{-i})$, and that the expected number of carries is
$\frac n2-\frac{1-b^{-n}}{2(b-1)}$. Evaluate it for two six-digit decimal numbers.
\end{enumerate}
\end{exercise}

\begin{exercise}[$\star\star$]\label{pb:laplace}
(Laplace, 1812.) A floor is tiled with squares of side $a$. A needle of length $\ell$ is thrown on it at random: its centre
is uniform, and its direction is uniform and independent of the centre.
\begin{enumerate}[label=(\alph*)]
\item For $\ell\le a$, show that the needle meets at least one line of the tiling with probability $\dfrac{4a\ell-\ell^2}{\pi a^2}$.
\item Show that for every $\ell>0$ the expected number of lines met is $\frac{4\ell}{\pi a}$, and deduce, for $\ell\le a$, the
probability that the needle meets two lines.
\end{enumerate}
\end{exercise}

\begin{exercise}[$\star\star$]\label{pb:detsquare}
Let $A$ be an $n\times n$ matrix whose entries are independent and uniform on $\{-1,1\}$.
\begin{enumerate}[label=(\alph*)]
\item Show that $\mathbb E[\det A]=0$ and $\mathbb E[(\det A)^2]=n!$.
\item Deduce that some $n\times n$ matrix with entries $\pm1$ has $|\det A|\ge\sqrt{n!}$.
\item Show that every real $n\times n$ matrix with entries in $[-1,1]$ satisfies $|\det A|\le n^{n/2}$ (Hadamard's
inequality), and compare the two bounds as $n\to\infty$.
\end{enumerate}
\end{exercise}

\section{Longer crossings}

The last four problems cross more than two borders. The first two are about the extreme cases of Problem~\ref{pb:detsquare}:
matrices of $\pm1$ whose rows are perfectly orthogonal, and polynomials with coefficients $\pm1$ that are as flat as possible on
the unit circle.

\begin{exercise}[$\star\star\star$]\label{pb:hadamard}
A \emph{Hadamard matrix} of order $n$ is an $n\times n$ matrix $H$ with entries $\pm1$ and $HH^{\mathsf T}=nI_n$.
\begin{enumerate}[label=(\alph*)]
\item Show that if a Hadamard matrix of order $n>2$ exists, then $4\mid n$.
\item Show that if $H$ is a Hadamard matrix of order $n$, then $\begin{pmatrix}H&H\\H&-H\end{pmatrix}$ is one of order $2n$.
\item (Paley, 1933.) Let $p\equiv3\pmod4$ be a prime, $\chi$ the Legendre symbol modulo $p$ (with $\chi(0)=0$), and $Q$ the
$p\times p$ matrix with entries $Q_{ab}=\chi(b-a)$, $a,b\in\mathbb F_p$. Show that $Q^{\mathsf T}=-Q$, that every row of $Q$
sums to $0$, and that $QQ^{\mathsf T}=pI_p-J_p$, where $J_p$ is the all-ones matrix. Deduce a Hadamard matrix of order $p+1$.
\item List the orders $n\le64$ obtained from (b) and (c).
\end{enumerate}
\end{exercise}

\begin{exercise}[$\star\star\star$]\label{pb:flat}
For signs $\varepsilon_0,\dots,\varepsilon_n\in\{-1,1\}$, let $P(z)=\sum_{k=0}^n\varepsilon_kz^k$ and
$\|P\|=\max_{|z|=1}|P(z)|$.
\begin{enumerate}[label=(\alph*)]
\item Show that $\|P\|\ge\sqrt{n+1}$ for every choice of signs.
\item Show that for every $n\ge2$ some choice of signs gives
\[
\|P\|\le2\sqrt{(n+1)\log\bigl(5n(n+1)\bigr)}+\frac\pi2.
\]
\end{enumerate}
\end{exercise}

\begin{exercise}[$\star\star$]\label{pb:rudin}
(Golay 1951, Shapiro 1952, Rudin 1959.) Let $P_0=Q_0=1$ and, for $k\ge0$,
\[
P_{k+1}(z)=P_k(z)+z^{2^k}Q_k(z),\qquad Q_{k+1}(z)=P_k(z)-z^{2^k}Q_k(z).
\]
\begin{enumerate}[label=(\alph*)]
\item Show that $P_k$ and $Q_k$ have degree $2^k-1$ and all their coefficients in $\{-1,1\}$.
\item Show that $|P_k(z)|^2+|Q_k(z)|^2=2^{k+1}$ for $|z|=1$, and deduce that $\|P_k\|\le\sqrt2\sqrt{2^k}$: within a factor
$\sqrt2$ of the lower bound of Problem~\ref{pb:flat}(a) (Figure~\ref{fig:rudin}).
\item Show that the coefficient of $z^m$ in $P_k$ (for $m<2^k$) is $(-1)^{e(m)}$, where $e(m)$ is the number of pairs of
adjacent $1$s in the binary expansion of $m$.
\end{enumerate}
\end{exercise}

\begin{figure}[ht]
\centering
\fig{ch01-rudin}
\caption{$|P(e^{i\theta})|$ for the Rudin--Shapiro polynomial $P_6$ (blue, $64$ coefficients) and for $64$ random signs (red).
The first never exceeds $\sqrt{128}\approx11.31$; the second reaches $15.3$. Both have mean square $64$.}
\label{fig:rudin}
\end{figure}

\begin{exercise}[$\star\star\star$]\label{pb:glnz}
(Finite groups of integer matrices.)
\begin{enumerate}[label=(\alph*)]
\item Show that $\cos\frac{2\pi}m$ is rational only for $m\in\{1,2,3,4,6\}$, and that every element of finite order of
$GL_2(\Z)$ has order $1$, $2$, $3$, $4$ or $6$.
\item Let $p\ge3$ be a prime and $A\in GL_n(\Z)$ with $A\equiv I_n\pmod p$ and $A\ne I_n$. Show that $A$ has infinite order.
Deduce that every finite subgroup of $GL_n(\Z)$ is isomorphic to a subgroup of $GL_n(\mathbb F_p)$, and that for $n=2$ its
order divides $48$.
\item Show that every finite subgroup of $GL_2(\Z)$ is cyclic or dihedral, of order at most $12$, and that $12$ is attained.
\end{enumerate}
\end{exercise}

\section{Hints}

\begin{hint}{pb:sicherman}
Encode a die by the polynomial $\sum_{\text{faces}}x^{\text{face}}$; the sum of independent dice corresponds to the product.
Factor $(x+x^2+\cdots+x^6)^2$ into irreducibles of $\Q[X]$, and evaluate the factors at $0$ and $1$.
\end{hint}
\begin{hint}{pb:loaded}
For $n=6$: the generating polynomials have degree $5$, and $1+x+\dots+x^{10}$ has no real root. For every $n$: compare the
coefficients of $x^0$, $x^{2n-2}$ and $x^{n-1}$ of the product, with the AM--GM inequality.
\end{hint}
\begin{hint}{pb:primepoly}
(a) If $P(a)=\pm p$, look at $P(a+kp)$ modulo $p$. (b) Evaluate $P$ at a multiple of $P(0)$ times the product of the primes
you already have.
\end{hint}
\begin{hint}{pb:fpfunctions}
$1-(x-a)^{p-1}$ is the indicator of $\{a\}$ on $\mathbb F_p$. For (c), use the power sums $\sum_xx^t$ modulo $p$.
\end{hint}
\begin{hint}{pb:fibmodp}
Binet's formula holds in any field where $X^2-X-1$ has two distinct roots; the Frobenius $x\mapsto x^p$ permutes them.
\end{hint}
\begin{hint}{pb:whichroot}
Vandermonde's identity $\sum_k\binom\alpha k\binom\beta{n-k}=\binom{\alpha+\beta}n$ is a polynomial identity in $\alpha,\beta$.
The two square roots of $8$ in $\Q_7$ differ modulo $7$, and those of $9$ in $\Q_2$ differ modulo $8$.
\end{hint}
\begin{hint}{pb:pascal2}
$(1+x)^{2^k}\equiv1+x^{2^k}$ in $\mathbb F_2[x]$.
\end{hint}
\begin{hint}{pb:millin}
Use $F_{m+n}=F_mF_{n+1}+F_{m-1}F_n$ and Cassini's identity $F_{m-1}F_{m+1}-F_m^2=(-1)^m$.
\end{hint}
\begin{hint}{pb:sylvester}
(b) $s_{k+1}-1=s_k(s_k-1)$. (c) Put $u_k=s_k-\frac12$; then $u_{k+1}=u_k^2+\frac14$. Study $2^{-k}\log u_k$.
\end{hint}
\begin{hint}{pb:smith}
$\gcd(i,j)=\sum_{d\mid i,\ d\mid j}\varphi(d)$. Write the matrix as $ADA^{\mathsf T}$ with $A$ triangular.
\end{hint}
\begin{hint}{pb:lambert}
Expand $\frac{x^n}{1-x^n}=\sum_{k\ge1}x^{kn}$, justify the regrouping, and recall the Dirichlet convolutions $\varphi*\mathbf1$,
$\mu*\mathbf1$ and $\lambda*\mathbf1$.
\end{hint}
\begin{hint}{pb:coprime52}
Every pair $(a,b)$ is $(ga',gb')$ with $g=\gcd(a,b)$ and $\gcd(a',b')=1$.
\end{hint}
\begin{hint}{pb:cyclewalk}
The indicator of $S\equiv0\pmod k$ is $\frac1k\sum_j\omega^{jS}$ with $\omega=e^{2i\pi/k}$.
\end{hint}
\begin{hint}{pb:fixmoments}
Burnside's lemma. An orbit of $k$-tuples is determined by which coordinates are equal.
\end{hint}
\begin{hint}{pb:matrices2}
Count conjugacy classes and their sizes with the orbit--stabiliser formula. A matrix that commutes with a non-scalar
$2\times2$ matrix $C$ is a polynomial in $C$.
\end{hint}
\begin{hint}{pb:carries}
$C_i=1$ exactly when $X_i+Y_i+C_{i-1}\ge b$, and $(X_i,Y_i)$ is independent of the past.
\end{hint}
\begin{hint}{pb:laplace}
Condition on the direction $\theta$: the projections of the needle on the axes are $\ell|\cos\theta|$ and $\ell|\sin\theta|$,
and the two coordinates of the centre are independent. For (b), use linearity and Buffon's needle.
\end{hint}
\begin{hint}{pb:detsquare}
Expand $\det A$ as a sum over permutations, then expand the square. For (c), Gram--Schmidt.
\end{hint}
\begin{hint}{pb:hadamard}
(a) Normalise the first row to all $1$s and count the columns according to the signs in rows $2$ and $3$. (c) For $d\ne0$,
$\sum_c\chi(c)\chi(c+d)=-1$. Put $S=\begin{pmatrix}0&\mathbf1^{\mathsf T}\\-\mathbf1&Q\end{pmatrix}$ and $H=I+S$.
\end{hint}
\begin{hint}{pb:flat}
(a) Parseval. (b) Random signs, Hoeffding's inequality for the real and imaginary parts at $M=n(n+1)$ equally spaced points,
and $|P'(z)|\le\frac{n(n+1)}2$ between them.
\end{hint}
\begin{hint}{pb:rudin}
The parallelogram law $|u+v|^2+|u-v|^2=2|u|^2+2|v|^2$. For (c), find and prove a formula for the coefficients of $Q_k$ too.
\end{hint}
\begin{hint}{pb:glnz}
(a) $2\cos\frac{2\pi}m$ is an algebraic integer, and so is the trace of a matrix of finite order. (b) Reduce to $A$ of prime
order $q$, write $A=I+p^kB$ with $p\nmid B$, and expand $A^q$. (c) Average a scalar product over the group.
\end{hint}

\section{Solutions}

\begin{solution}{pb:sicherman}
\idea A die is a multiset of faces, and the law of a sum of independent variables is the product of the generating functions
(Tome~3, Chapter~3). So the question is how $(x+\dots+x^6)^2$ factors. Unique factorisation in $\Q[X]$ leaves only finitely
many candidates, and the conditions ``six faces'' and ``positive faces'' are the values at $1$ and at $0$.

\writeup To a die with faces $a_1,\dots,a_6$ attach $A(x)=\sum_ix^{a_i}$; its coefficients are non-negative integers,
$A(1)=6$ and, since the faces are positive, $A(0)=0$. If $B$ is the polynomial of the second die, the number of ways to obtain
the sum $s$ is the coefficient of $x^s$ in $A(x)B(x)$, so the condition is
\[
A(x)B(x)=(x+x^2+\dots+x^6)^2=x^2(1+x)^2(1+x+x^2)^2(1-x+x^2)^2,
\]
because $x+\dots+x^6=x\frac{x^6-1}{x-1}=x(1+x)(1+x+x^2)(1-x+x^2)$. The four factors $x$, $1+x$, $1+x+x^2$, $1-x+x^2$ are
irreducible in $\Q[X]$: the quadratics have negative discriminant, hence no rational root. Since $\Q[X]$ has unique
factorisation (Tome~0, Chapter~3), $A=\lambda x^a(1+x)^b(1+x+x^2)^c(1-x+x^2)^d$ with $\lambda\in\Q^*$ and $0\le a,b,c,d\le2$,
and $B=\lambda'x^{2-a}(1+x)^{2-b}(1+x+x^2)^{2-c}(1-x+x^2)^{2-d}$ with $\lambda\lambda'=1$. The leading coefficients $\lambda$
and $\lambda'$ of $A$ and $B$ are positive integers, so $\lambda=\lambda'=1$.

Now evaluate. At $x=1$ the factors take the values $1,2,3,1$, so $A(1)=2^b3^c=6$ forces $b=c=1$. At $x=0$, $A(0)=0$ forces
$a\ge1$ and $B(0)=0$ forces $2-a\ge1$, so $a=1$. Only $d\in\{0,1,2\}$ is free. For $d=1$ both dice are ordinary. For $d=0$,
$A=x(1+x)(1+x+x^2)=x+2x^2+2x^3+x^4$, the die $1$, $2$, $2$, $3$, $3$, $4$, and then $d'=2$ gives
$B=x(1+x)(1+x+x^2)(1-x+x^2)^2=x+x^3+x^4+x^5+x^6+x^8$, the die $1$, $3$, $4$, $5$, $6$, $8$; $d=2$ is the same pair in the other order. Both
polynomials have non-negative coefficients, so they are dice. There are exactly two solutions: the ordinary pair and
Sicherman's pair.
\tools Tome~3, Chapter~3 (generating functions of sums); Tome~0, Chapter~3 (unique factorisation in $\Q[X]$).
\leads The same method classifies ``Sicherman dice'' with $n$ faces: there are as many pairs as ways of splitting the
cyclotomic factors $\Phi_d$, $d\mid n$, $d>1$, of $x\frac{x^n-1}{x-1}$ between two dice so that both have non-negative
coefficients (Gallian and Rusin, 1979).
\end{solution}

\begin{solution}{pb:loaded}
\idea Again the generating polynomials multiply. A real polynomial of odd degree has a real root, while the polynomial of the
uniform law on $2,\dots,12$ has none: that settles $n=6$. For general $n$, only three coefficients matter: the extreme ones are
products $p_1q_1$ and $p_nq_n$, and the middle one contains $p_1q_n+p_nq_1$, which AM--GM bounds from below.

\writeup Let $A(x)=\sum_{i=1}^np_ix^{i-1}$ and $B(x)=\sum_jq_jx^{j-1}$. If the sum were uniform on $\{2,\dots,2n\}$, then
\[
A(x)B(x)=\frac1{2n-1}\bigl(1+x+\dots+x^{2n-2}\bigr).
\]
\emph{First proof, $n=6$.} The right-hand side has degree $10$, so $A$ and $B$, of degree at most $5$, both have degree $5$.
A real polynomial of odd degree has a real root (intermediate value theorem), so $AB$ would have a real root. But
$1+x+\dots+x^{10}=\frac{x^{11}-1}{x-1}$ for $x\ne1$, whose roots are the $11$-th roots of unity other than $1$, none of them
real, and its value at $1$ is $11$.

\emph{Second proof, every $n\ge2$.} Comparing the coefficients of $x^0$ and $x^{2n-2}$ gives $p_1q_1=p_nq_n=\frac1{2n-1}$; in
particular all four numbers are positive. The coefficient of $x^{n-1}$ is $\sum_{i+j=n+1}p_iq_j\ge p_1q_n+p_nq_1$, and by
AM--GM
\[
p_1q_n+p_nq_1\ge2\sqrt{p_1q_np_nq_1}=2\sqrt{p_1q_1\cdot p_nq_n}=\frac2{2n-1}>\frac1{2n-1},
\]
a contradiction with the coefficient $\frac1{2n-1}$ of $x^{n-1}$ on the right.
\tools Tome~3, Chapter~3 (generating functions); Tome~0, Chapter~4 (AM--GM) and Chapter~3 (roots of polynomials).
\end{solution}

\begin{solution}{pb:primepoly}
\idea A polynomial with integer coefficients respects congruences: $P(a+kp)\equiv P(a)\pmod p$. So once $P(a)=\pm p$, all the
values $P(a+kp)$ are multiples of $p$; if they were all prime they would all be $\pm p$, too many times for a non-constant
polynomial. For (b), Euclid's trick: build a value of $P$ that is $\equiv P(0)$ modulo a big product, then divide out
$P(0)$.

\writeup (a) Suppose that $|P(n)|$ is prime for every $n\ge N$. Let $a=N$ and $p=|P(a)|$. For $k\ge0$,
$a+kp\equiv a\pmod p$, so $P(a+kp)\equiv P(a)\equiv0\pmod p$ (Tome~1, Chapter~3). Since $|P(a+kp)|$ is a prime divisible by
$p$, it equals $p$. So one of the polynomials $P-p$, $P+p$ has infinitely many roots, hence is zero (Tome~0, Chapter~3), and
$P$ is constant: a contradiction.

(b) Let $c=P(0)$. If $c=0$, then $n\mid P(n)$ for every $n$, so every prime $p$ divides $P(p)$. Let $c\ne0$, and suppose that
only the primes $p_1,\dots,p_r$ divide values of $P$. Put $m=cp_1\cdots p_r$ and write $P(X)=c+XR(X)$ with $R\in\Z[X]$. For
every integer $t$,
\[
P(mt)=c+mtR(mt)=c\bigl(1+p_1\cdots p_r\,tR(mt)\bigr).
\]
The polynomial $t\mapsto1+p_1\cdots p_r\,tR(mt)$ is not constant ($P$ is not), so it takes the values $\pm1$ only finitely
often, and for some $t$ the integer $1+p_1\cdots p_rtR(mt)$ has a prime factor $q$. It divides $P(mt)$, so $q$ is some $p_i$;
but $1+p_1\cdots p_r\,tR(mt)\equiv1\pmod{p_i}$. This contradiction shows that infinitely many primes divide values of $P$.
\tools Tome~1, Chapters~2--3 (primes, congruences); Tome~0, Chapter~3 (a non-zero polynomial has finitely many roots).
\leads Part (b) is Schur's theorem (1912). Which primes divide values of $P$ is a deep question: for $P=X^2+1$ they are $2$
and the primes $\equiv1\pmod4$, and in general the answer is governed by the Galois group of $P$ (Chebotarev's density
theorem, Tome~12). No polynomial of degree $\ge2$ is known to take infinitely many prime values; for $X^2+1$ this is one of
Landau's four problems (1912), still open.
\end{solution}

\begin{solution}{pb:fpfunctions}
\idea The exponent $p-1$ is special in $\mathbb F_p$: by Fermat, $y^{p-1}$ is $1$ for $y\ne0$ and $0$ for $y=0$. So
$1-(x-a)^{p-1}$ is the indicator of the point $a$, and any function is a combination of indicators: this is Lagrange
interpolation in its most economical form. Its top coefficient is read off at once, and a bijection has the same multiset of
values as the identity.

\writeup (a) Let $\delta_a(X)=1-(X-a)^{p-1}$. For $x\in\mathbb F_p$, $(x-a)^{p-1}=1$ if $x\ne a$ (Fermat's little theorem,
Tome~1, Chapter~3) and $0$ if $x=a$, so $\delta_a(x)=[x=a]$. Hence $P=\sum_af(a)\delta_a$, of degree at most $p-1$, satisfies
$P(x)=f(x)$ for all $x$. If two polynomials of degree at most $p-1$ define the same map, their difference has $p$ roots and
degree at most $p-1$, so it is zero (Tome~0, Chapter~3; the proof works over any field). So $P$ is unique. (As a check: there
are $p^p$ maps and $p^p$ polynomials of degree at most $p-1$.)

(b) The coefficient of $X^{p-1}$ in $\delta_a$ is $-1$, so that of $P$ is $-\sum_af(a)$.

(c) If $f$ is a bijection, $\sum_af(a)=\sum_xx=\frac{p(p-1)}2$, which is $0$ in $\mathbb F_p$ for $p$ odd; by (b), $\deg P\le
p-2$. For $1\le t\le p-2$, the map $x\mapsto f(x)^t$ has $\sum_af(a)^t=\sum_xx^t=0$ in $\mathbb F_p$, because the power sum
$\sum_{x=1}^{p-1}x^t$ is $\equiv0\pmod p$ when $(p-1)\nmid t$ (Tome~1, Chapter~5). By (b) again its polynomial has degree at
most $p-2$. For $f(x)=x^3$ on $\mathbb F_5$ ($f$ is a bijection since $\gcd(3,4)=1$): the polynomials of $f$, $f^2$, $f^3$ are
obtained by reducing $X^3$, $X^6$, $X^9$ with $x^5=x$, which gives $X^3$, $X^2$ and $X$, all of degree at most $3$.
\tools Tome~1, Chapters~3 and~5 (Fermat, power sums modulo $p$); Tome~0, Chapter~3 (Lagrange interpolation, roots).
\leads A polynomial that permutes $\mathbb F_q$ is a \emph{permutation polynomial}. Hermite's criterion (1863) says that the
conditions of (c), together with ``$f$ has exactly one root'', characterise them. They are used in cryptography and in the
construction of codes (Lidl and Niederreiter, \emph{Finite Fields}, Chapter~7).
\end{solution}

\begin{solution}{pb:fibmodp}
\idea Binet's formula $F_n=\frac{\varphi^n-\psi^n}{\varphi-\psi}$ is an identity between sequences satisfying the same
recurrence, so it holds in every field where $X^2-X-1$ has two distinct roots, for instance in a field with $p$ or $p^2$
elements. In such a field the Frobenius $x\mapsto x^p$ is an automorphism; it either fixes the roots (when $5$ is a square
modulo $p$) or swaps them. Both cases give a power of $\varphi$ equal to $1$ or to $-1$.

\writeup The discriminant of $X^2-X-1$ is $5\ne0$ in $\mathbb F_p$. If $\left(\frac5p\right)=1$ its roots lie in $K=\mathbb
F_p$; otherwise it is irreducible over $\mathbb F_p$ and its roots lie in $K=\mathbb F_p[X]/(X^2-X-1)$, a field with $p^2$
elements (Tome~2, Chapter~9). Call the roots $\varphi\ne\psi$; then $\varphi+\psi=1$ and $\varphi\psi=-1$. The sequences
$(F_n\bmod p)$ and $\bigl(\frac{\varphi^n-\psi^n}{\varphi-\psi}\bigr)$ both satisfy $u_{n+2}=u_{n+1}+u_n$ in $K$, with
$u_0=0$, $u_1=1$; so they coincide. The map $x\mapsto x^p$ is a field homomorphism of $K$ fixing $\mathbb F_p$ (Tome~2,
Chapter~9), so it maps a root of $X^2-X-1$ to a root.

\emph{Case $\left(\frac5p\right)=1$.} Then $\varphi,\psi\in\mathbb F_p^*$, so $\varphi^{p-1}=\psi^{p-1}=1$ by Fermat. Hence
$F_{n+p-1}=F_n$ in $\mathbb F_p$ for all $n$, and $F_{p-1}=F_0=0$.

\emph{Case $\left(\frac5p\right)=-1$.} Then $\varphi\notin\mathbb F_p$. The elements of $K$ fixed by $x\mapsto x^p$ are the
roots of $X^p-X$, at most $p$ of them, and the $p$ elements of $\mathbb F_p$ are roots; so $\varphi^p\ne\varphi$, and
$\varphi^p=\psi$. Likewise $\psi^p=\varphi$. So $\varphi^{p+1}=\varphi\psi=-1=\psi^{p+1}$, which gives
$F_{p+1}=\frac{-1-(-1)}{\varphi-\psi}=0$, and $\varphi^{2(p+1)}=\psi^{2(p+1)}=1$, so $F_{n+2(p+1)}=F_n$ for all $n$.

This proves (a). For (b): the set of $N\in\Z$ such that $F_{n+N}\equiv F_n$ for all $n\ge0$ (extend the sequence to negative
indices by the recurrence) is a subgroup of $\Z$, generated by the period; it contains $p-1$, respectively $2(p+1)$, so the
period divides it. Finally, by quadratic reciprocity $\left(\frac5p\right)=\left(\frac p5\right)$ (Tome~1, Chapter~6), and the
squares modulo $5$ are $\pm1$; this gives the two cases as stated.

(c) For $p=7\equiv2\pmod5$, the period divides $16$; the sequence modulo $7$ is $0$, $1$, $1$, $2$, $3$, $5$, $1$, $6$, $0$, $6$, $6$, $5$, $4$, $2$, $6$, $1$, $0$, $1$, \dots,
of period $16$. For $p=11\equiv1\pmod5$ it divides $10$: $0$, $1$, $1$, $2$, $3$, $5$, $8$, $2$, $10$, $1$, $0$, $1$, \dots, of period $10$.
\tools Tome~2, Chapter~9 (finite fields, Frobenius); Tome~1, Chapter~6 (quadratic reciprocity); Tome~0, Chapter~7 (linear
recurrences, Binet).
\leads The period $\pi(p)$ is called the Pisano period. Wall asked in 1960 whether $\pi(p^2)\ne\pi(p)$ for every prime $p$; a
prime with $\pi(p^2)=\pi(p)$ would be a \emph{Wall--Sun--Sun prime}, and none is known (they are linked to the first case of
Fermat's last theorem). The notebook computes $\pi(p)$ for all $p<10^5$.
\end{solution}

\begin{solution}{pb:whichroot}
\idea The binomial series $\sum\binom{1/2}nx^n$ is ``$\sqrt{1+x}$'' because its square is $1+x$ as a formal power series. In
$\Q_7$, the powers $7^n$ are small, so the series converges; in $\R$ it diverges, since $|x|>1$. Which root it picks is decided
by the first term and the size of the others, and the surprise is that in $\Q_2$ the series for $\sqrt9$ lands on $-3$.

\writeup (a) For $n\ge1$, the numerator of $\binom{1/2}n$ is
\[
\frac12\cdot\Bigl(-\frac12\Bigr)\Bigl(-\frac32\Bigr)\cdots\Bigl(-\frac{2n-3}2\Bigr)=(-1)^{n-1}\frac{1\cdot3\cdots(2n-3)}{2^n},
\]
and $1\cdot3\cdots(2n-3)=\frac{(2n-2)!}{2^{n-1}(n-1)!}$. So
\[
\binom{1/2}n=(-1)^{n-1}\frac{(2n-2)!}{2^{2n-1}\,n!\,(n-1)!}=(-1)^{n-1}\frac{C_{n-1}}{2^{2n-1}},
\]
since $C_{n-1}=\frac{(2n-2)!}{n!(n-1)!}$, an integer (Tome~3, Chapter~1).

\emph{The square of the series.} For non-negative integers $\alpha,\beta$, $\sum_{k=0}^n\binom\alpha k\binom\beta{n-k}=
\binom{\alpha+\beta}n$: choose $n$ people among $\alpha$ women and $\beta$ men. Both sides are polynomials in $(\alpha,\beta)$
that agree on $\N^2$, so they are equal for all $\alpha,\beta$ (a polynomial in two variables vanishing on $\N^2$ is zero:
apply Tome~0, Chapter~3 in each variable). With $\alpha=\beta=\frac12$: $\sum_k\binom{1/2}k\binom{1/2}{n-k}=\binom1n$, which is
$1$ for $n=0,1$ and $0$ for $n\ge2$.

\emph{Products in $\Q_p$.} If $a_n\to0$ and $b_n\to0$ in $\Q_p$, then $\sum a_n$ and $\sum b_n$ converge (Tome~1,
Chapter~12), and so does $\sum c_n$ with $c_n=\sum_{k\le n}a_kb_{n-k}$, to the product of the sums. Indeed, let $\beta=\sup|b_n|$
and $\alpha=\sup|a_n|$. Then
$\bigl(\sum_{k\le N}a_k\bigr)\bigl(\sum_{l\le N}b_l\bigr)-\sum_{n\le N}c_n=\sum a_kb_l$ over $k,l\le N$ with $k+l>N$; in each
term $k>\frac N2$ or $l>\frac N2$, so by the ultrametric inequality the difference has absolute value at most
$\max\bigl(\beta\sup_{k>N/2}|a_k|,\alpha\sup_{l>N/2}|b_l|\bigr)\to0$.

(b) The term $\binom{1/2}n7^n$ has $7$-adic absolute value $7^{-n}$ at most, since $C_{n-1}$ is an integer and $2^{2n-1}$ is
a unit in $\Z_7$. So the series converges to some $s\in\Z_7$, and by the two facts above $s^2=1+7=8$. Every term with $n\ge1$
lies in $7\Z_7$, so $s\equiv1\pmod7$. The square roots of $8$ in $\Q_7$ are $\pm s$, and $-s\equiv-1\not\equiv1$: the series
converges to \emph{the} square root of $8$ that is $\equiv1\pmod7$.

(c) Now $\binom{1/2}n8^n=(-1)^{n-1}C_{n-1}2^{3n-(2n-1)}=(-1)^{n-1}C_{n-1}2^{n+1}$, whose $2$-adic absolute value is at most
$2^{-n-1}$. The series converges in $\Z_2$ to some $s$ with $s^2=9$, so $s=\pm3$. Modulo $8$: the terms with $n\ge2$ are
divisible by $2^{n+1}\ge8$, and the first two are $1$ and $\binom{1/2}18=4$. So $s\equiv5\pmod8$, while $3\not\equiv5$ and
$-3\equiv5\pmod8$. Hence $s=-3$.
\tools Tome~1, Chapter~12 (convergence in $\Q_p$); Tome~0, Chapters~1 and~3 (binomial identities, polynomials); Tome~3,
Chapter~1 (Catalan numbers).
\leads The same series can converge in $\R$ and in several $\Q_p$ to different square roots, a first hint that a rational
number ``lives'' in all completions of $\Q$ at once. This point of view leads to the adeles (Tome~7) and to the local--global
principle (Tome~12).
\end{solution}

\begin{solution}{pb:pascal2}
\idea Row $n$ of Pascal's triangle is the polynomial $(1+x)^n$. Modulo $2$, squaring is additive (Frobenius), so
$(1+x)^{2^k}=1+x^{2^k}$, and multiplying by it copies a row twice, shifted. That is the self-similar picture; the count follows
by peeling off binary digits one at a time.

\writeup (a) In $\mathbb F_2[x]$, $(u+v)^2=u^2+v^2$, so by induction $(1+x)^{2^k}=1+x^{2^k}$ (Tome~2, Chapter~9). Hence, for
$0\le m<2^k$,
\[
(1+x)^{2^k+m}=(1+x^{2^k})(1+x)^m=(1+x)^m+x^{2^k}(1+x)^m\quad\text{in }\mathbb F_2[x],
\]
and comparing the coefficients of $x^j$ gives the claim. As $(1+x)^m$ has degree $m<2^k$, the two copies do not overlap: row
$2^k+m$ is row $m$, then a gap of zeros, then row $m$ again. This is the pattern of Figure~\ref{fig:pascal}: the triangle of
rows $0$ to $2^{k+1}-1$ consists of three copies of the triangle of rows $0$ to $2^k-1$.

(b) Let $N(n)$ be the number of odd entries of row $n$. By (a), $N(2^k+m)=2N(m)$ for $0\le m<2^k$, and $N(0)=1$. Removing the
leading binary digit of $n$ halves $N$ and decreases $s$ by one, so $N(n)=2^{s(n)}$ by induction on $n$. (This is also
Lucas's theorem, Tome~1, Chapter~3: $\binom nj$ is odd iff every binary digit of $j$ is at most the corresponding digit of
$n$, which gives $2^{s(n)}$ choices of $j$.)

(c) Write $n<2^k$ with $k$ binary digits $n_0,\dots,n_{k-1}$. Then $2^{s(n)}=\prod_i2^{n_i}$, and summing over all digit
choices gives $\sum_{n<2^k}2^{s(n)}=\prod_{i<k}(1+2)=3^k$. The rows $0$ to $2^k-1$ contain $\sum_{n<2^k}(n+1)=2^{k-1}(2^k+1)$
entries, so the proportion of odd ones is $\frac{3^k}{2^{k-1}(2^k+1)}\le2\bigl(\frac34\bigr)^k\to0$.
\tools Tome~2, Chapter~9 (Frobenius in characteristic $p$); Tome~1, Chapter~3 (Lucas's theorem); Tome~0, Chapter~1.
\leads The limit shape of Figure~\ref{fig:pascal}, rescaled, is the Sierpiński triangle, of dimension $\frac{\log3}{\log2}$
(Tome~6 defines the dimension). Modulo a prime $p$ the same argument gives $\prod(n_i+1)$ non-zero entries in row $n$ (Fine,
1947).
\end{solution}

\begin{solution}{pb:millin}
\idea A sum with denominators $F_{2^n}$, growing doubly exponentially, smells of telescoping: one looks for a quantity $T_N$
with $T_N-T_{N+1}=\frac1{F_{2^{N+1}}}$. Computing the first partial sums $2$, $\frac73$, $\frac{50}{21}$ and comparing them
with $3$ suggests $T_N=\frac{F_{2^N-1}}{F_{2^N}}$, and the doubling formulas prove it.

\writeup (a) The addition formula $F_{m+n}=F_mF_{n+1}+F_{m-1}F_n$ (Tome~0, Chapter~1) with $n=m$ gives
$F_{2m}=F_mF_{m+1}+F_{m-1}F_m$, and with $n=m-1$ gives $F_{2m-1}=F_m^2+F_{m-1}^2$.

(b) Let $m\ge2$ be even. Then
\[
\frac{F_{m-1}}{F_m}-\frac{F_{2m-1}}{F_{2m}}=\frac{F_{m-1}F_m(F_{m-1}+F_{m+1})-F_m(F_m^2+F_{m-1}^2)}{F_mF_{2m}}
=\frac{F_{m-1}F_{m+1}-F_m^2}{F_{2m}}=\frac{(-1)^m}{F_{2m}}=\frac1{F_{2m}},
\]
by Cassini's identity (Tome~0, Chapter~7). With $m=2^N$, $N\ge1$, this says $\bigl(3-\frac{F_{2^{N+1}-1}}{F_{2^{N+1}}}\bigr)-
\bigl(3-\frac{F_{2^N-1}}{F_{2^N}}\bigr)=\frac1{F_{2^{N+1}}}$, and the formula for $N=1$ is
$\frac1{F_1}+\frac1{F_2}=2=3-\frac{F_1}{F_2}$; induction gives the claim. Finally $\frac{F_{n-1}}{F_n}\to\frac1\varphi=
\frac{\sqrt5-1}2$ (Tome~4, Chapter~2), so the sum is $3-\frac{\sqrt5-1}2=\frac{7-\sqrt5}2=2.3819660\ldots$
\tools Tome~0, Chapters~1 and~7 (Fibonacci identities, Cassini); Tome~4, Chapter~2 (Binet's formula).
\leads This is Millin's series (1974). The sum $\sum_{n\ge1}\frac1{F_n}=3.3598856\ldots$ has no known closed form; André-Jeannin
proved in 1989 that it is irrational.
\end{solution}

\begin{solution}{pb:sylvester}
\idea For (a), something must decrease: here it is the numerator of the remainder. For (b), the recursion $s_{k+1}-1=
s_k(s_k-1)$ turns $\frac1{s_k}$ into a difference, exactly as $\frac1{n(n+1)}=\frac1n-\frac1{n+1}$. For (c), a sequence that
squares at each step is $E^{2^k}$ to first order; taking logarithms turns the small correction $\frac14$ into a convergent
series.

\writeup (a) Let $n=\lceil\frac ba\rceil\ge2$, so that $\frac1n\le\frac ab<\frac1{n-1}$. The remainder is
$\frac ab-\frac1n=\frac{an-b}{bn}$, and $0\le an-b<a$ because $a(n-1)<b$. So the numerators (before reduction) decrease
strictly and the process stops. The remainder is $<\frac1{n-1}-\frac1n=\frac1{n(n-1)}\le\frac1n$, so the next unit fraction
is smaller than $\frac1n$: the fractions are distinct. For $\frac4{13}$: $n=4$, remainder $\frac3{52}$; then $n=18$,
remainder $\frac3{52}-\frac1{18}=\frac1{468}$. So $\frac4{13}=\frac14+\frac1{18}+\frac1{468}$.

(b) Since $s_{k+1}-1=s_k^2-s_k=s_k(s_k-1)$, we get $\frac1{s_k-1}-\frac1{s_{k+1}-1}=\frac{s_k-1}{s_k(s_k-1)}=\frac1{s_k}$, and the
sum telescopes: $\sum_{k=0}^N\frac1{s_k}=\frac1{s_0-1}-\frac1{s_{N+1}-1}=1-\frac1{s_{N+1}-1}$. Iterating the same identity,
$s_{k+1}-1=s_ks_{k-1}\cdots s_0(s_0-1)=s_0s_1\cdots s_k$, so $s_{k+1}\equiv1\pmod{s_j}$ for every $j\le k$: the $s_k$ are
pairwise coprime. (Their prime factors are therefore all distinct, which proves again that there are infinitely many primes,
as the Fermat numbers do in Tome~1, Chapter~2.)

(c) Let $u_k=s_k-\frac12$. Then $u_{k+1}=s_k^2-s_k+\frac12=u_k^2+\frac14$, and $u_0=\frac32$, so $(u_k)$ increases to infinity.
Let $c_k=2^{-k}\log u_k$. Then
\[
c_{k+1}-c_k=2^{-k-1}\log\Bigl(1+\frac1{4u_k^2}\Bigr)\in\Bigl]0,\frac{2^{-k-1}}{4u_k^2}\Bigr],
\]
a summable positive sequence; so $c_k$ increases to a limit $c$. Put $E=e^{c/2}$, so $E^{2^{k+1}}=e^{2^kc}$. Then
\[
2^kc-\log u_k=2^k\sum_{j\ge k}(c_{j+1}-c_j)=\sum_{j\ge k}2^{k-j-1}\log\Bigl(1+\frac1{4u_j^2}\Bigr),
\]
which lies in $\bigl]0,\frac1{4u_k^2}\bigr]$, because $u_j\ge u_k$ and $\sum_{j\ge k}2^{k-j-1}=1$. Hence
$u_k<E^{2^{k+1}}\le u_ke^{1/(4u_k^2)}$. Since $e^t-1\le te^t$ and $u_k\ge\frac32$,
\[
u_k\bigl(e^{1/(4u_k^2)}-1\bigr)\le\frac{e^{1/9}}{4u_k}\le\frac{1.12}6<\frac12.
\]
So $s_k-\frac12<E^{2^{k+1}}<s_k$, that is $s_k<E^{2^{k+1}}+\frac12<s_k+\frac12$, and $s_k=\bigl\lfloor E^{2^{k+1}}+\frac12
\bigr\rfloor$. Numerically $E=1.2640847353\ldots$
\tools Tome~1, Chapters~1--2 (coprimality, infinitely many primes); Tome~0, Chapter~5 (monotone sequences, series);
Tome~4, Chapter~11 (asymptotics of fast recursions).
\leads Aho and Sloane (1973) studied such doubly exponential recursions. Erdős and Graham asked many questions on Egyptian
fractions; one of them, that every $\frac4n$ is a sum of three unit fractions (Erdős--Straus, 1948), is still open.
\end{solution}

\begin{solution}{pb:smith}
\idea A matrix whose entries are a function of $\gcd(i,j)$ is a sum over common divisors, which is a matrix product: one
factor records ``$d$ divides $i$'', the other ``$d$ divides $j$''. The divisibility matrix is triangular with ones on the
diagonal, so the determinant is the product of the diagonal weights.

\writeup (b) Let $A$ be the $n\times n$ matrix with $A_{id}=1$ if $d\mid i$ and $0$ otherwise, and $D=\operatorname{diag}
(f(1),\dots,f(n))$. Then
\[
(ADA^{\mathsf T})_{ij}=\sum_{d=1}^nA_{id}f(d)A_{jd}=\sum_{d\mid i,\ d\mid j}f(d)=\sum_{d\mid\gcd(i,j)}f(d)=g(\gcd(i,j)),
\]
since the common divisors of $i$ and $j$ are the divisors of $\gcd(i,j)$, all at most $n$. The matrix $A$ is lower triangular
($d\mid i$ implies $d\le i$) with $A_{ii}=1$, so $\det A=1$ and $\det\bigl(g(\gcd(i,j))\bigr)=\det D=f(1)\cdots f(n)$ (Tome~0,
Chapter~7).

(a) is the case $f=\varphi$, since $\sum_{d\mid m}\varphi(d)=m$ (Tome~1, Chapter~4).

For $\gcd(i,j)^2$: we need $f$ with $f*\mathbf1=\mathrm{id}^2$, that is $f=\mu*\mathrm{id}^2$ by Möbius inversion, the Jordan
function $J_2(m)=\sum_{d\mid m}\mu(d)\bigl(\frac md\bigr)^2=m^2\prod_{p\mid m}\bigl(1-\frac1{p^2}\bigr)$. So
$\det\bigl(\gcd(i,j)^2\bigr)=\prod_{m=1}^nJ_2(m)$; for $n=4$ this is $1\cdot3\cdot8\cdot12=288$.
\tools Tome~1, Chapter~4 (Dirichlet convolution, Möbius inversion); Tome~0, Chapter~7 (determinants of triangular matrices).
\leads H.~J.~S.~Smith found this in 1875. The same factorisation shows the matrix $(\gcd(i,j))$ is positive definite, and
the asymptotic behaviour of its eigenvalues is still studied (Lindqvist and Seip, 1998).
\end{solution}

\begin{solution}{pb:lambert}
\idea $\frac{x^n}{1-x^n}=x^n+x^{2n}+x^{3n}+\cdots$ has a term $x^m$ exactly when $n\mid m$. So the Lambert series collects, in
front of $x^m$, the sum of $a_d$ over the divisors $d$ of $m$: it is the power series of $a*\mathbf1$. Choosing $a$ with a
known $a*\mathbf1$ gives closed forms.

\writeup (a) For $|x|<1$ the family $\bigl(a_nx^{kn}\bigr)_{n,k\ge1}$ is absolutely summable:
$\sum_n|a_n|\sum_k|x|^{kn}=\sum_n|a_n|\frac{|x|^n}{1-|x|^n}\le\frac1{1-|x|}\sum_nCn^K|x|^n<\infty$. So it may be summed in any
order (Tome~0, Chapter~5): summing over $k$ first gives the Lambert series, and grouping the terms with $kn=m$ gives
$\sum_m\bigl(\sum_{d\mid m}a_d\bigr)x^m$.

(b) Since $\varphi*\mathbf1=\mathrm{id}$, $\mu*\mathbf1=[m=1]$ (Tome~1, Chapter~4):
\[
\sum_n\varphi(n)\frac{x^n}{1-x^n}=\sum_mmx^m=\frac x{(1-x)^2},\qquad\sum_n\mu(n)\frac{x^n}{1-x^n}=x.
\]
For $\lambda$: it is completely multiplicative, so $\lambda*\mathbf1$ is multiplicative, and on a prime power
$\sum_{j=0}^a\lambda(p^j)=1-1+\dots\pm1$ is $1$ if $a$ is even and $0$ if $a$ is odd. So $(\lambda*\mathbf1)(m)=1$ if $m$ is a
square and $0$ otherwise, and $\sum_n\lambda(n)\frac{x^n}{1-x^n}=\sum_{j\ge1}x^{j^2}$. (Similarly $\Lambda*\mathbf1=\log$ gives
$\sum_m(\log m)x^m$.)

(c) At $x=\frac12$, $\frac{x^n}{1-x^n}=\frac1{2^n-1}$, so the two sums are $\frac{1/2}{1/4}=2$ and $\frac12$.
\tools Tome~1, Chapter~4 (convolution identities); Tome~0, Chapter~5 and Tome~4, Chapter~2 (summable families, power series).
\leads The Lambert series of $\lambda$ is a theta function; Lambert series of divisor functions are Eisenstein series, the
first modular forms (Tome~5, Chapter~12, and Tome~12).
\end{solution}

\begin{solution}{pb:coprime52}
\idea Pull out the gcd: summing over all pairs is summing over $g$ and over coprime pairs, and the gcd contributes a factor
$\zeta(2s)$. So the coprime sum is $\frac{\zeta(s)^2}{\zeta(2s)}$, and at $s=2$ the values of $\zeta(2)$ and $\zeta(4)$ do the
rest.

\writeup Let $s>1$. All the series have positive terms, so they may be rearranged freely. Every pair $(a,b)$ of positive
integers is uniquely $(ga',gb')$ with $g=\gcd(a,b)$ and $\gcd(a',b')=1$. Hence
\[
\zeta(s)^2=\sum_{a,b\ge1}\frac1{(ab)^s}=\sum_{g\ge1}\frac1{g^{2s}}\sum_{\gcd(a',b')=1}\frac1{(a'b')^s}=\zeta(2s)\,S(s),
\]
so $S(s)=\frac{\zeta(s)^2}{\zeta(2s)}$. With $\zeta(2)=\frac{\pi^2}6$ and $\zeta(4)=\frac{\pi^4}{90}$ (Tome~4, Chapter~3),
$S(2)=\frac{\pi^4}{36}\cdot\frac{90}{\pi^4}=\frac52$.

Grouping by the product $n=a'b'$, $S(s)=\sum_nc_nn^{-s}$ where $c_n$ is the number of ordered coprime pairs $(a',b')$ with
$a'b'=n$. Coprimality forces each prime power $p^e\parallel n$ to go entirely to $a'$ or to $b'$, so $c_n=2^{\omega(n)}$, where
$\omega(n)$ is the number of distinct prime factors of $n$. So $\sum_n2^{\omega(n)}n^{-s}=\frac{\zeta(s)^2}{\zeta(2s)}$.
\tools Tome~1, Chapters~1 and~4 (gcd, multiplicative functions); Tome~4, Chapter~3 ($\zeta(2)$, $\zeta(4)$).
\end{solution}

\begin{solution}{pb:cyclewalk}
\idea The position after $n$ steps is a sum $S_n=X_1+\dots+X_n$ of independent $\pm1$, read modulo $k$. To pick out the event
$S_n\equiv0$ one uses the roots of unity filter, and the expectation of $\omega^{jS_n}$ factorises by independence.

\writeup (a) Let $\omega=e^{2i\pi/k}$. For an integer $s$, $\frac1k\sum_{j=0}^{k-1}\omega^{js}$ is $1$ if $k\mid s$ and $0$
otherwise (Tome~0, Chapter~2). So
\[
q_n=\mathbb E\Bigl[\frac1k\sum_j\omega^{jS_n}\Bigr]=\frac1k\sum_j\prod_{i=1}^n\mathbb E\bigl[\omega^{jX_i}\bigr]
=\frac1k\sum_j\Bigl(\frac{\omega^j+\omega^{-j}}2\Bigr)^n=\frac1k\sum_j\cos^n\frac{2\pi j}k,
\]
using independence (Tome~3, Chapter~2).

(b) For $k=3$ the cosines are $1,-\frac12,-\frac12$, so $q_n=\frac13\bigl(1+2(-\frac12)^n\bigr)$. For $k=4$ they are
$1,0,-1,0$, so $q_n=\frac14\bigl(1+(-1)^n\bigr)$ for $n\ge1$ (and $q_0=1$). For $k=6$ they are $1,\frac12,-\frac12,-1,-\frac12,
\frac12$, so $q_n=\frac16\bigl(1+2\cdot2^{-n}+2(-\frac12)^n+(-1)^n\bigr)$.

(c) $\cos\frac{2\pi j}k=\pm1$ only for $j=0$ and, if $k$ is even, $j=\frac k2$. The other terms tend to $0$. So $q_n\to\frac1k$
if $k$ is odd; if $k$ is even, $q_n\to\frac2k$ along even $n$ and $q_n=0$ for odd $n$ (an even polygon is bipartite).
\tools Tome~0, Chapter~2 (roots of unity filter); Tome~3, Chapter~2 (independence) and Chapter~8 (periodic chains).
\end{solution}

\begin{solution}{pb:fixmoments}
\idea $F(\sigma)^k$ counts the $k$-tuples of fixed points of $\sigma$, that is the $k$-tuples fixed by $\sigma$ acting on
$\{1,\dots,n\}^k$. Averaging a number of fixed points over a group is exactly what Burnside's lemma turns into a number of
orbits. The orbits are equality patterns, that is set partitions, and the Bell numbers are the moments of the Poisson law.

\writeup (a) A tuple $x$ is fixed by $\sigma$ iff $\sigma(x_i)=x_i$ for every $i$, iff every $x_i$ is a fixed point of
$\sigma$; there are $F(\sigma)^k$ such tuples. By Burnside's lemma (Tome~2, Chapter~3), the number of orbits is
$\frac1{n!}\sum_\sigma F(\sigma)^k=\mathbb E[F^k]$.

(b) To a tuple $x$ attach the partition $\pi(x)$ of $\{1,\dots,k\}$ whose blocks are the sets $\{i:x_i=v\}$, $v$ a value taken
by $x$. A permutation does not change it, since $\sigma(x_i)=\sigma(x_j)$ iff $x_i=x_j$. Conversely, if $\pi(x)=\pi(y)$, the map
$x_i\mapsto y_i$ is a well-defined bijection from the set of values of $x$ onto that of $y$, and extends to a permutation
$\sigma$ of $\{1,\dots,n\}$ (both complements have the same size), with $\sigma\cdot x=y$. So the orbits correspond to the
partitions of $\{1,\dots,k\}$ into at most $n$ blocks (a tuple takes at most $n$ values, and every such partition occurs).
There are $\stirlingii kj$ partitions into $j$ blocks (Tome~4, Chapter~1), so $\mathbb E[F^k]=\sum_{j\le n}\stirlingii kj$,
which is $\sum_j\stirlingii kj=B_k$ when $k\le n$, since then no partition has more than $n$ blocks.

(c) If $N$ is Poisson of parameter $1$, $\mathbb E[N^k]=e^{-1}\sum_{m\ge0}\frac{m^k}{m!}=B_k$ by Dobiński's formula (Tome~4,
Chapter~2). So $\mathbb E[F^k]=\mathbb E[N^k]$ for $1\le k\le n$; for $k=n+1$ they differ, since $\mathbb E[F^{n+1}]=B_{n+1}-1$
(the partition into $n+1$ singletons is missing).
\tools Tome~2, Chapter~3 (Burnside's lemma); Tome~4, Chapters~1--2 (Stirling and Bell numbers, Dobiński); Tome~3, Chapter~2.
\leads Moments that agree up to order $n$ imply a total variation distance between $F$ and Poisson(1) of order
$\frac{2^n}{n!}$, super-exponentially small (Diaconis and Pitman; Arratia and Tavaré, 1992). The notebook compares the two laws.
\end{solution}

\begin{solution}{pb:matrices2}
\idea Group the matrices by conjugacy: $GL_2(\mathbb F_p)$ acts on $M_2(\mathbb F_p)$ by conjugation, the class of a matrix is
determined by its characteristic polynomial and whether it is scalar, and orbit--stabiliser gives the size of each class
once the centraliser is known. For a non-scalar $C$, the centraliser is the invertible part of $\mathbb F_p[C]$, a ring with
$p^2$ elements: $\mathbb F_p\times\mathbb F_p$, $\mathbb F_{p^2}$ or $\mathbb F_p[\varepsilon]/(\varepsilon^2)$.

\writeup Let $G=GL_2(\mathbb F_p)$, of order $(p^2-1)(p^2-p)=p(p-1)^2(p+1)$ (Tome~2, Chapter~1).

\emph{Centralisers.} Let $C$ be non-scalar. Some vector $v$ is not an eigenvector, so $(v,Cv)$ is a basis. If $AC=CA$,
write $Av=\alpha v+\beta Cv$; then $B=A-\alpha I-\beta C$ commutes with $C$ and kills $v$, hence kills $Cv$ ($BCv=CBv=0$), so
$B=0$. So the centraliser of $C$ in $M_2(\mathbb F_p)$ is $\mathbb F_p[C]=\{\alpha I+\beta C\}$, which has $p^2$ elements,
and its centraliser in $G$ is the group of invertible elements of $\mathbb F_p[C]$.

\emph{The classes.} The $p$ scalar matrices are alone in their classes. A non-scalar $M$ with characteristic polynomial
$\chi$ is conjugate to the companion matrix of $\chi$ (in the basis $(v,Mv)$ above), so the non-scalar matrices with a given
$\chi$ form one class, of size $|G|/u(\chi)$, where $u(\chi)$ is the number of invertible elements of $\mathbb F_p[C]
\cong\mathbb F_p[X]/(\chi)$.
\begin{itemize}
\item $\chi=(X-\lambda)(X-\mu)$, $\lambda\ne\mu$: $\mathbb F_p[X]/(\chi)\cong\mathbb F_p\times\mathbb F_p$ (Chinese remainders),
$u=(p-1)^2$, class size $p(p+1)$. There are $\frac{p(p-1)}2$ such $\chi$: in all $\frac{p^2(p^2-1)}2$ matrices.
\item $\chi$ irreducible: $\mathbb F_p[X]/(\chi)$ is a field with $p^2$ elements, $u=p^2-1$, class size $p(p-1)$. There are
$\frac{p^2-p}2$ monic irreducible quadratics (Tome~2, Chapter~9): in all $\frac{p^2(p-1)^2}2$ matrices.
\item $\chi=(X-\lambda)^2$, $M$ not scalar: $\mathbb F_p[X]/((X-\lambda)^2)$ has $p(p-1)$ invertible elements (those with
non-zero value at $\lambda$), class size $p^2-1$, and $p$ choices of $\lambda$: in all $p(p^2-1)$ matrices.
\end{itemize}
As a check, $p+\frac{p^2(p^2-1)}2+\frac{p^2(p-1)^2}2+p^3-p=p^4$.

(b) $M$ is diagonalisable over $\mathbb F_p$ iff it is scalar or has two distinct eigenvalues in $\mathbb F_p$ (a non-scalar
matrix with a double eigenvalue is not, since $M-\lambda I\ne0$ is nilpotent). So
$P(\text{diagonalisable})=\frac{p+p^2(p^2-1)/2}{p^4}\to\frac12$ and $P(\chi\text{ irreducible})=\frac{(p-1)^2}{2p^2}\to\frac12$.
\tools Tome~2, Chapter~3 (orbit--stabiliser), Chapter~6 (Chinese remainders in $K[X]$), Chapter~9 (counting irreducible
polynomials); Tome~0, Chapter~7 (diagonalisation).
\leads The two limits are the probabilities that a random permutation of the two roots is the identity or a transposition.
This is no accident: Chapter~5 shows that the degrees of the factors of a random polynomial over $\mathbb F_q$ follow the cycle
type of a random permutation, as $q\to\infty$.
\end{solution}

\begin{solution}{pb:carries}
\idea Whether position $i$ produces a carry depends only on its two digits and on the carry coming in, and the two digits
are fresh randomness. That is the definition of a Markov chain. Counting the pairs of digits that overflow gives the
transition probabilities; the chain is symmetric, so its stationary law is $(\frac12,\frac12)$, and the second eigenvalue
measures how fast it is reached.

\writeup (a) Let $X_i,Y_i$ be the digits at position $i$. Then $C_i=1$ iff $X_i+Y_i+C_{i-1}\ge b$. Since $(X_i,Y_i)$ is
independent of $(C_0,\dots,C_{i-1})$, which is a function of the earlier digits, $P(C_i=1\mid C_0,\dots,C_{i-1})$ depends only
on $C_{i-1}$: the chain is Markov (Tome~3, Chapter~8). The number of pairs $(x,y)$ of digits with $x+y\ge b$ is
$\sum_{x=1}^{b-1}x=\frac{b(b-1)}2$, and with $x+y\ge b-1$ it is $\frac{b(b+1)}2$. So
\[
P=\frac1{2b}\begin{pmatrix}b+1&b-1\\b-1&b+1\end{pmatrix},
\]
with rows indexed by the carry in ($0$, then $1$). It is symmetric and stochastic, with eigenvalues $1$ (eigenvector
$(1,1)$) and $\frac{(b+1)-(b-1)}{2b}=\frac1b$ (eigenvector $(1,-1)$).

(b) Let $q_i=P(C_i=1)$. Then $q_i=(1-q_{i-1})\frac{b-1}{2b}+q_{i-1}\frac{b+1}{2b}=\frac{b-1}{2b}+\frac{q_{i-1}}b$, so
$q_i-\frac12=\frac1b\bigl(q_{i-1}-\frac12\bigr)$, and $q_0=0$ gives $q_i=\frac12(1-b^{-i})$. The expected number of carries
is $\sum_{i=1}^nq_i=\frac n2-\frac12\sum_{i=1}^nb^{-i}=\frac n2-\frac{1-b^{-n}}{2(b-1)}$. For $b=10$ and $n=6$, this is
$3-\frac{1-10^{-6}}{18}=2.9444445$.
\tools Tome~3, Chapter~8 (Markov chains, eigenvalues); Tome~1, Chapter~1 (base-$b$ digits).
\leads Holte (1997) found that adding $k$ numbers gives a chain with eigenvalues $1,b^{-1},\dots,b^{-(k-1)}$ and a stationary law
given by Eulerian numbers; Diaconis and Fulman (2009) showed that it is the chain of descents of riffle shuffles
(Tome~3, Chapter~8).
\end{solution}

\begin{solution}{pb:laplace}
\idea Fix the direction first. The needle then occupies a rectangle of sides $\ell|\cos\theta|$ and $\ell|\sin\theta|$, and it
avoids the vertical lines iff the horizontal coordinate of its centre stays far enough from them; the two coordinates are
independent, so the probability of avoiding everything is a product. Averaging over $\theta$ is an elementary integral. For
the expected number of crossings, linearity reduces everything to Buffon's single family of lines.

\writeup (a) By periodicity we may take the centre $(x,y)$ uniform in $[0,a]^2$, and the angle $\theta$ with the horizontal
uniform in $[0,\pi[$, all independent (Tome~3, Chapter~5). Given $\theta$, the needle meets no vertical line iff
$\frac\ell2|\cos\theta|\le x\le a-\frac\ell2|\cos\theta|$, which has probability $1-\frac\ell a|\cos\theta|$ since $\ell\le a$;
and likewise for horizontal lines with $\sin\theta$. By independence of $x$ and $y$,
\[
P(\text{no line}\mid\theta)=\Bigl(1-\frac\ell a|\cos\theta|\Bigr)\Bigl(1-\frac\ell a|\sin\theta|\Bigr).
\]
By symmetry we may average over $[0,\frac\pi2]$, where $\int\cos=\int\sin=1$ and $\int\sin\theta\cos\theta\dd\theta=\frac12$:
\[
P(\text{no line})=\frac2\pi\Bigl(\frac\pi2-\frac{2\ell}a+\frac{\ell^2}{2a^2}\Bigr)=1-\frac{4\ell}{\pi a}+\frac{\ell^2}{\pi a^2},
\]
so the needle meets a line with probability $\frac{4a\ell-\ell^2}{\pi a^2}$.

(b) Let $N_v$ and $N_h$ be the numbers of vertical and horizontal lines met. For the family of vertical lines alone, Buffon's
computation gives $\mathbb E[N_v]=\frac{2\ell}{\pi a}$ for every $\ell$ (Tome~3, Chapter~1: given $\theta$, the expected number of
crossings is $\frac{\ell|\cos\theta|}a$, whose mean is $\frac{2\ell}{\pi a}$). The same holds for $N_h$, so $\mathbb E[N]=
\frac{4\ell}{\pi a}$. For $\ell\le a$, $N_v,N_h\le1$, so $N\in\{0,1,2\}$ and $\mathbb E[N]=P(N\ge1)+P(N=2)$. Hence
\[
P(N=2)=\frac{4\ell}{\pi a}-\frac{4a\ell-\ell^2}{\pi a^2}=\frac{\ell^2}{\pi a^2}.
\]
\tools Tome~3, Chapter~1 (Buffon's needle) and Chapter~5 (uniform laws, independence); Tome~0, Chapter~6 (integrals).
\leads Note that $P(N=2)=\frac{\ell^2}{\pi a^2}$ differs from $P(N_v=1)P(N_h=1)=\bigl(\frac{2\ell}{\pi a}\bigr)^2$: the two
crossings are dependent, through the common angle $\theta$. Integral geometry (Santaló, \emph{Integral Geometry and Geometric Probability}, 1976)
computes such probabilities for any shapes.
\end{solution}

\begin{solution}{pb:detsquare}
\idea The determinant is a signed sum of $n!$ products. Squaring it creates pairs of products, and the expectation kills every
pair in which some entry appears an odd number of times. Only the $n!$ diagonal pairs survive, each contributing $1$. A random
variable that has mean square $n!$ must sometimes be at least $\sqrt{n!}$ in absolute value.

\writeup (a) $\det A=\sum_\sigma\operatorname{sgn}(\sigma)\prod_ia_{i\sigma(i)}$. Each product has mean $0$, its factors being
independent with mean $0$. Next,
\[
(\det A)^2=\sum_{\sigma,\tau}\operatorname{sgn}(\sigma\tau)\prod_ia_{i\sigma(i)}a_{i\tau(i)}.
\]
If $\sigma=\tau$ the product is $\prod_ia_{i\sigma(i)}^2=1$. If $\sigma\ne\tau$, pick $i$ with $\sigma(i)\ne\tau(i)$: the entry
$a_{i\sigma(i)}$ appears exactly once in the product (row $i$ contributes $a_{i\sigma(i)}a_{i\tau(i)}$, and the other rows
contribute other rows' entries), so the product is $a_{i\sigma(i)}$ times an independent variable, and has mean $0$. Hence
$\mathbb E[(\det A)^2]=n!$ (Tome~3, Chapter~2).

(b) If $|\det A|<\sqrt{n!}$ for every $\pm1$ matrix, then $\mathbb E[(\det A)^2]<n!$.

(c) Let $r_1,\dots,r_n$ be the rows. If they are dependent, $\det A=0$. Otherwise, Gram--Schmidt gives
$r_i^*=r_i-\sum_{j<i}c_{ij}r_j^*$ with $r_1^*,\dots,r_n^*$ orthogonal (Tome~0, Chapter~7). Subtracting from a row a combination
of the previous rows does not change the determinant, so $\det A=\det A^*$ with $A^*$ the matrix of rows $r_i^*$, and
$(\det A^*)^2=\det(A^*A^{*\mathsf T})=\prod_i\|r_i^*\|^2$. As $r_i^*$ is the orthogonal projection of $r_i$ on the orthogonal
of the previous rows, $\|r_i^*\|\le\|r_i\|\le\sqrt n$. So $|\det A|\le n^{n/2}$. By Stirling, $\frac{\sqrt{n!}}{n^{n/2}}\sim
(2\pi n)^{1/4}e^{-n/2}$: the typical size of a random determinant is exponentially smaller than the maximum.
\tools Tome~3, Chapter~2 (expectation, independence); Tome~0, Chapter~7 (determinants, Gram--Schmidt) and Chapters~5--6 (Stirling's formula).
\leads Tao and Vu (2006) proved that $|\det A|=\sqrt{n!}\,e^{O(\sqrt{n\log n})}$ with probability tending to $1$, and Nguyen and
Vu (2014) that $\log|\det A|$ is asymptotically normal. Equality in Hadamard's bound is Problem~\ref{pb:hadamard}.
\end{solution}

\begin{solution}{pb:hadamard}
\idea A Hadamard matrix is a set of $n$ pairwise orthogonal $\pm1$ vectors. Orthogonality of $\pm1$ vectors means they agree
in exactly half of the coordinates, and three rows force a split into quarters. For a construction, one needs $\pm1$
vectors that are ``uncorrelated'', and the quadratic residues modulo $p$ are the classical source: shifting the Legendre
symbol by $d$ leaves it almost orthogonal to itself (Jacobsthal's sum). A border of ones and a diagonal fix the ``almost''.

\writeup (a) Multiplying a column by $-1$ keeps the Hadamard property ($HH^{\mathsf T}=\sum_c(\text{column }c)(\text{column
}c)^{\mathsf T}$), so we may assume row $1$ is all ones. Let $n\ge3$, and let $\alpha,\beta,\gamma,\delta$ be the numbers of
columns where rows $2$ and $3$ carry $(+,+)$, $(+,-)$, $(-,+)$, $(-,-)$. Orthogonality of rows $1,2,3$ gives
\[
\alpha+\beta-\gamma-\delta=0,\qquad\alpha-\beta+\gamma-\delta=0,\qquad\alpha-\beta-\gamma+\delta=0,
\]
and $\alpha+\beta+\gamma+\delta=n$. Adding all four, $4\alpha=n$; similarly $\beta=\gamma=\delta=\frac n4$. So $4\mid n$.

(b) By block multiplication,
\[
\begin{pmatrix}H&H\\H&-H\end{pmatrix}\begin{pmatrix}H^{\mathsf T}&H^{\mathsf T}\\H^{\mathsf T}&-H^{\mathsf T}\end{pmatrix}
=\begin{pmatrix}2HH^{\mathsf T}&0\\0&2HH^{\mathsf T}\end{pmatrix}=2nI_{2n}.
\]

(c) Since $p\equiv3\pmod4$, $\chi(-1)=-1$ (Tome~1, Chapter~6), so $Q_{ba}=\chi(a-b)=-\chi(b-a)=-Q_{ab}$. Each row sums to
$\sum_b\chi(b-a)=\sum_c\chi(c)=0$, since there are as many residues as non-residues. For the diagonal of $QQ^{\mathsf T}$:
$\sum_c\chi(c-a)^2=p-1$. Off the diagonal, with $d=a-b\ne0$ and $c$ replaced by $c+b$:
\[
(QQ^{\mathsf T})_{ab}=\sum_c\chi(c-a)\chi(c-b)=\sum_c\chi(c)\chi(c+d).
\]
For $c\ne0$, $\chi(c)\chi(c+d)=\chi(c^2)\chi(1+dc^{-1})=\chi(1+dc^{-1})$, and as $c$ runs over $\mathbb F_p^*$, $1+dc^{-1}$
runs over $\mathbb F_p\setminus\{1\}$. So the sum is $\sum_{u\ne1}\chi(u)=-\chi(1)=-1$. Hence $QQ^{\mathsf T}=pI-J$.

Now let $S=\begin{pmatrix}0&\mathbf1^{\mathsf T}\\-\mathbf1&Q\end{pmatrix}$, of order $p+1$, and $H=I+S$. The diagonal of
$Q$ is $\chi(0)=0$, so $H$ has entries $\pm1$. Next $S^{\mathsf T}=-S$, and
\[
SS^{\mathsf T}=\begin{pmatrix}0&\mathbf1^{\mathsf T}\\-\mathbf1&Q\end{pmatrix}\begin{pmatrix}0&-\mathbf1^{\mathsf T}\\
\mathbf1&Q^{\mathsf T}\end{pmatrix}=\begin{pmatrix}p&(Q\mathbf1)^{\mathsf T}\\Q\mathbf1&J+QQ^{\mathsf T}\end{pmatrix}=pI_{p+1},
\]
since $Q\mathbf1=0$. Therefore $HH^{\mathsf T}=(I+S)(I-S)=I-S^2=I+SS^{\mathsf T}=(p+1)I_{p+1}$.

(d) Paley's construction gives the orders $4$, $8$, $12$, $20$, $24$, $32$, $44$, $48$, $60$ (from $p=3,7,11,19,23,31,43,47,59$), Sylvester's doubling
of $1$ gives $1$, $2$, $4$, $8$, $16$, $32$, $64$, and doubling the Paley orders gives $24,40,48,64$ and then $48$ again. So up to $64$ we obtain
every multiple of $4$ except $28$, $36$, $52$ and $56$. (Those exist too, by other constructions of Paley and Williamson.)
\tools Tome~1, Chapter~6 (Legendre symbol, $\chi(-1)$, residue counts); Tome~0, Chapter~7 (block matrices); Tome~2,
Chapter~9 (the field $\mathbb F_p$).
\leads The \emph{Hadamard conjecture} says that a Hadamard matrix exists for every multiple of $4$; it is open, and the
smallest unknown order is $668$ (the order $428$ was found by Kharaghani and Tayfeh-Rezaie in 2005). In computer science the
Sylvester matrices are the Walsh--Hadamard transform, and their rows form the Hadamard code (Tome~C4).
\end{solution}

\begin{solution}{pb:flat}
\idea Parseval says the mean of $|P|^2$ on the circle is the sum of the squares of the coefficients, so the maximum is at
least $\sqrt{n+1}$. For an upper bound, choose the signs at random: at each fixed point of the circle $P$ is a sum of
independent bounded terms, so Hoeffding makes a large value exponentially unlikely, and a union bound handles finitely many
points. Between the points, $P$ cannot move much because its derivative is bounded.

\writeup (a) For $z=e^{i\theta}$, we have $|P(z)|^2=\sum_{k,l}\varepsilon_k\varepsilon_le^{i(k-l)\theta}$, and the integral of
$e^{im\theta}$ over a period is $0$ for $m\ne0$. So
\[
\frac1{2\pi}\int_0^{2\pi}|P(e^{i\theta})|^2\dd\theta=\sum_k\varepsilon_k^2=n+1
\]
(Tome~4, Chapter~6), and $\|P\|^2$, being at least this mean, is at least $n+1$.

(b) Let the $\varepsilon_k$ be independent fair signs. Fix $\theta$. The variables $\varepsilon_k\cos k\theta$ are independent,
centred, with values in $[-|\cos k\theta|,|\cos k\theta|]$, so Hoeffding's inequality (Tome~3, Chapter~4) gives
\[
P\bigl(|\operatorname{Re}P(e^{i\theta})|\ge t\bigr)\le2\exp\Bigl(-\frac{2t^2}{\sum_k4\cos^2k\theta}\Bigr)\le2e^{-t^2/(2(n+1))},
\]
and the same for the imaginary part. Let $M=n(n+1)$ and $\theta_j=\frac{2\pi j}M$, $0\le j<M$. If $|\operatorname{Re}P|<t$ and
$|\operatorname{Im}P|<t$ then $|P|<\sqrt2t$, so
\[
P\bigl(\exists j:|P(e^{i\theta_j})|\ge\sqrt2t\bigr)\le4Me^{-t^2/(2(n+1))}.
\]
With $t^2=2(n+1)\log(5M)$ the right-hand side is $\frac45<1$, so some choice of signs has $|P(e^{i\theta_j})|<\sqrt2t=
2\sqrt{(n+1)\log(5M)}$ for every $j$. For this choice and any $\theta$, let $\theta_j$ be the nearest grid point, so
$|\theta-\theta_j|\le\frac\pi M$. Since $|P'(z)|\le\sum_kk=\frac{n(n+1)}2$ on the circle,
\[
|P(e^{i\theta})-P(e^{i\theta_j})|\le|\theta-\theta_j|\cdot\frac{n(n+1)}2\le\frac\pi M\cdot\frac M2=\frac\pi2.
\]
So $\|P\|\le2\sqrt{(n+1)\log(5n(n+1))}+\frac\pi2$.
\tools Tome~3, Chapter~4 (Hoeffding's inequality, union bound); Tome~4, Chapter~6 (Parseval); Tome~0, Chapter~2.
\leads Salem and Zygmund (1954) showed that for random signs $\|P\|$ is of order $\sqrt{n\log n}$, so (b) is sharp for random
choices; Problem~\ref{pb:rudin} removes the logarithm with a clever choice. Littlewood asked in 1966 for signs with
$\delta\sqrt n\le|P(z)|\le\Delta\sqrt n$ on the whole circle; Balister, Bollobás, Morris, Sahasrabudhe and Tiba proved in 2020
that such ``flat Littlewood polynomials'' exist.
\end{solution}

\begin{solution}{pb:rudin}
\idea The recursion is built so that $(P_{k+1},Q_{k+1})$ is $(P_k+w,P_k-w)$ with $|w|=|Q_k|$: the parallelogram law then
doubles $|P|^2+|Q|^2$ exactly, at every point of the circle. For the coefficients, the recursion concatenates coefficient lists,
and one guesses the pattern from the first cases: $P_2=1+z+z^2-z^3$, $P_3=1+z+z^2-z^3+z^4+z^5-z^6+z^7$.

\writeup (a) By induction: $P_k$ and $Q_k$ have degree $<2^k$, and $z^{2^k}Q_k$ only has terms of degree between $2^k$ and
$2^{k+1}-1$; so the coefficient list of $P_{k+1}$ is that of $P_k$ followed by that of $Q_k$, and of $Q_{k+1}$ that of $P_k$
followed by that of $-Q_k$. All are $\pm1$, and the top coefficient is non-zero.

(b) For $|z|=1$, let $u=P_k(z)$ and $v=z^{2^k}Q_k(z)$, so $|v|=|Q_k(z)|$. Then $|P_{k+1}|^2+|Q_{k+1}|^2=|u+v|^2+|u-v|^2
=2|u|^2+2|v|^2=2\bigl(|P_k|^2+|Q_k|^2\bigr)$, and $|P_0|^2+|Q_0|^2=2$. So the sum is $2^{k+1}$, and $|P_k(z)|\le\sqrt{2^{k+1}}$.
Problem~\ref{pb:flat}(a) with $n+1=2^k$ gives $\|P_k\|\ge\sqrt{2^k}$.

(c) For $m<2^k$ let $t_k(m)$ be the binary digit of $m$ at position $k-1$ (the top position). We prove by induction on $k$
that for $m<2^k$,
\[
[z^m]P_k=(-1)^{e(m)},\qquad[z^m]Q_k=(-1)^{e(m)+t_k(m)}.
\]
For $k=0$, $m=0$ and both coefficients are $1$. Assume it for $k$, and let $m<2^{k+1}$. If $m<2^k$, then $[z^m]P_{k+1}=
[z^m]Q_{k+1}=[z^m]P_k=(-1)^{e(m)}$, and $t_{k+1}(m)=0$: correct. If $m=2^k+m'$ with $m'<2^k$, then $e(m)=e(m')+t_k(m')$ (the new
top digit $1$ forms a new pair exactly when the digit below it is $1$), and $t_{k+1}(m)=1$. So
$[z^m]P_{k+1}=[z^{m'}]Q_k=(-1)^{e(m')+t_k(m')}=(-1)^{e(m)}$ and $[z^m]Q_{k+1}=-[z^{m'}]Q_k=(-1)^{e(m)+1}$: correct.
\tools Tome~0, Chapter~2 (modulus of complex numbers); Tome~4, Chapter~6 (Parseval, through Problem~\ref{pb:flat});
Tome~1, Chapter~1 (binary expansions).
\leads Golay invented such ``complementary pairs'' in 1951 for infrared spectrometry; they are used in radar and Wi-Fi
signals. The Rudin--Shapiro sequence $(-1)^{e(m)}$ is $2$-automatic, and Mauduit and Rivat (2015) proved that it takes the
values $\pm1$ equally often on the primes.
\end{solution}

\begin{solution}{pb:glnz}
\idea An integer matrix of finite order has eigenvalues that are roots of unity, and its trace is an integer: that leaves only
a few angles. Reducing modulo $p$ is a homomorphism, and the question is whether it can kill an element of finite order; the
$p$-adic size of $A-I$ shows it cannot, because raising $I+p^kB$ to a power $q$ multiplies the ``error'' $p^kB$ by $q$ to
first order and only adds higher powers of $p$. Finally, a finite group of real matrices preserves some scalar product,
so it is a group of rotations and reflections in disguise.

\writeup (a) $2\cos\frac{2\pi}m=\zeta+\zeta^{-1}$ with $\zeta=e^{2i\pi/m}$, a sum of algebraic integers, hence an algebraic
integer (Tome~2, Chapter~11). If it is rational, it is an integer (Tome~1, Chapter~1), in $[-2,2]$, so $\cos\frac{2\pi}m\in
\{0,\pm\frac12,\pm1\}$ and $m\in\{1,2,3,4,6\}$.

Let $A\in GL_2(\Z)$ with $A^m=I$. Its minimal polynomial divides $X^m-1$, which has distinct roots in $\C$, so $A$ is
diagonalisable over $\C$ (Tome~0, Chapter~7), with eigenvalues $\lambda_1,\lambda_2$ that are roots of unity, and the order of
$A$ is the lcm of their orders. The characteristic polynomial has integer coefficients. If $\lambda_1$ is real, so is
$\lambda_2$, both are $\pm1$, and the order is $1$ or $2$. Otherwise $\lambda_2=\overline{\lambda_1}$, and
$\operatorname{tr}A=2\operatorname{Re}\lambda_1=2\cos\frac{2\pi a}{d}$ with $d$ the order of $\lambda_1$ and $\gcd(a,d)=1$.
This is an integer in $]-2,2[$, so the cosine is $0$ or $\pm\frac12$, and $d\in\{3,4,6\}$ (the angle $\frac{2\pi a}d$ is
$\pm\frac\pi2$, $\pm\frac\pi3$ or $\pm\frac{2\pi}3$ modulo $2\pi$). So the order of $A$ is in $\{1,2,3,4,6\}$.

(b) Suppose $A$ has finite order $m>1$, and let $q$ be a prime dividing $m$; then $A'=A^{m/q}$ has order $q$ and still
$A'\equiv I\pmod p$. Rename it $A$. Write $A=I+p^kB$ with $k\ge1$ and $B\in M_n(\Z)$ not divisible by $p$. Then
\[
I=A^q=I+qp^kB+\sum_{j=2}^q\binom qjp^{jk}B^j,\qquad\text{so}\qquad qp^kB=-\sum_{j=2}^q\binom qjp^{jk}B^j.
\]
If $q\ne p$: the right-hand side is divisible by $p^{2k}$, while $qp^kB$ is divisible by $p^k$ but not by $p^{k+1}$ (some
entry of $B$ is prime to $p$, and so is $q$). As $2k\ge k+1$, this is impossible. If $q=p$: for $2\le j<p$, $p\mid\binom pj$, so
the term is divisible by $p^{1+jk}$, with $1+jk\ge k+2$; the last term is divisible by $p^{pk}$, and $pk\ge k+2$ because
$p\ge3$. So the right-hand side is divisible by $p^{k+2}$, and $p^{k+1}B$ is not: impossible again.

So the reduction map $GL_n(\Z)\to GL_n(\mathbb F_p)$, a group homomorphism, is injective on every finite subgroup $G$, whose
order therefore divides $|GL_n(\mathbb F_p)|$ (Lagrange). For $n=2$, $|GL_2(\mathbb F_3)|=8\cdot6=48$, so $|G|$ divides $48$.
(For every $p\ge5$, $|GL_2(\mathbb F_p)|=p(p-1)^2(p+1)$ is a multiple of $48$, so other primes add nothing.)

(c) Let $G\subset GL_2(\Z)$ be finite. For $x,y\in\R^2$ put $\langle x,y\rangle_G=\sum_{g\in G}(gx)\cdot(gy)$. This is a scalar
product (positive definite, as a sum of positive definite forms), and it is $G$-invariant: $\langle hx,hy\rangle_G=\langle x,
y\rangle_G$ for $h\in G$, by reindexing $g\mapsto gh$. In an orthonormal basis for $\langle\,,\rangle_G$ the elements of $G$
become orthogonal matrices, so $G$ is isomorphic to a finite subgroup of $O(2)$. Let $G^+$ be its subgroup of rotations
(determinant $1$), of index $1$ or $2$. Rotations form a group isomorphic to the unit circle of $\C$, whose finite subgroups are
cyclic (Tome~2, Chapter~1), so $G^+$ is cyclic of some order $m$, generated by a rotation of order $m$, and by (a)
$m\in\{1,2,3,4,6\}$. If $G\ne G^+$, let $s\in G$ be a reflection; then $G=G^+\cup sG^+$ and $srs^{-1}=r^{-1}$ for every rotation
$r$, so $G$ is dihedral of order $2m$. Hence $|G|\le12$. The order $12$ is attained: $A=\begin{psmallmatrix}0&-1\\1&1
\end{psmallmatrix}$ (characteristic polynomial $X^2-X+1$, order $6$) and $B=\begin{psmallmatrix}0&1\\1&0\end{psmallmatrix}$
satisfy $BAB^{-1}=A^{-1}$ and generate a dihedral group of order $12$, the symmetries of the hexagonal lattice.
\tools Tome~2, Chapters~1 and~11 (finite subgroups of $\C^*$, algebraic integers); Tome~1, Chapters~1 and~3 (rational
algebraic integers, congruences, binomial coefficients modulo $p$); Tome~0, Chapter~7 (diagonalisation, scalar products).
\leads This is the \emph{crystallographic restriction}: a lattice in the plane has no symmetry of order $5$ or $7$, which is
why quasicrystals (Shechtman, Nobel prize 2011) were a shock. Minkowski (1887) used (b) for all primes at once to find the
largest finite subgroups of $GL_n(\Q)$, and Bieberbach's theorems (1911) classify crystallographic groups in all dimensions.
\end{solution}

\section*{Map of the chapter}

\begin{center}
\small
\adjustbox{max width=\linewidth}{\begin{tabular}{lll}
Problem & Topic & Tomes used\\\hline
\ref{pb:sicherman} & Sicherman's dice & 3, 0\\
\ref{pb:loaded} & loaded dice & 3, 0\\
\ref{pb:primepoly} & prime values of polynomials & 1, 0\\
\ref{pb:fpfunctions} & functions on $\mathbb F_p$ & 1, 0, 2\\
\ref{pb:fibmodp} & Fibonacci modulo $p$ & 2, 1, 0\\
\ref{pb:whichroot} & $p$-adic binomial series & 1, 0, 3\\
\ref{pb:pascal2} & Pascal's triangle modulo $2$ & 2, 1, 0\\
\ref{pb:millin} & Millin's series & 0, 4\\
\ref{pb:sylvester} & Egyptian fractions, Sylvester & 1, 0, 4\\
\ref{pb:smith} & Smith's determinant & 1, 0\\
\ref{pb:lambert} & Lambert series & 1, 0, 4\\
\ref{pb:coprime52} & a coprime double sum & 1, 4\\
\ref{pb:cyclewalk} & walk on a polygon & 0, 3\\
\ref{pb:fixmoments} & moments of fixed points & 2, 3, 4\\
\ref{pb:matrices2} & random $2\times2$ matrices over $\mathbb F_p$ & 2, 0\\
\ref{pb:carries} & carries & 3, 1\\
\ref{pb:laplace} & Laplace's needle & 3, 0\\
\ref{pb:detsquare} & random determinants & 3, 0\\
\ref{pb:hadamard} & Hadamard matrices & 1, 0, 2\\
\ref{pb:flat} & flat polynomials & 3, 4, 0\\
\ref{pb:rudin} & Rudin--Shapiro & 0, 4, 1\\
\ref{pb:glnz} & finite subgroups of $GL_2(\Z)$ & 2, 1, 0\\
\end{tabular}}
\end{center}

\section*{Sources and further reading}

\begin{itemize}
\item M.~Gardner, Mathematical Games, \emph{Scientific American} 238 (February 1978), and J.~A.~Gallian and D.~J.~Rusin,
Cyclotomic polynomials and nonstandard dice, \emph{Discrete Mathematics} 27 (1979), 245--259.
\item I.~Schur, Über die Existenz unendlich vieler Primzahlen in einigen speziellen arithmetischen Progressionen,
\emph{Sitzungsberichte der Berliner Mathematischen Gesellschaft} 11 (1912), 40--50.
\item R.~Lidl and H.~Niederreiter, \emph{Finite Fields}, 2nd ed., Cambridge, 1997, Chapter~7 (permutation polynomials, Hermite's
criterion).
\item D.~D.~Wall, Fibonacci series modulo $m$, \emph{American Mathematical Monthly} 67 (1960), 525--532.
\item N.~J.~Fine, Binomial coefficients modulo a prime, \emph{American Mathematical Monthly} 54 (1947), 589--592.
\item A.~V.~Aho and N.~J.~A.~Sloane, Some doubly exponential sequences, \emph{Fibonacci Quarterly} 11 (1973), 429--437.
\item H.~J.~S.~Smith, On the value of a certain arithmetical determinant, \emph{Proceedings of the London Mathematical Society}
7 (1875), 208--212.
\item J.~M.~Holte, Carries, combinatorics, and an amazing matrix, \emph{American Mathematical Monthly} 104 (1997), 138--149;
P.~Diaconis and J.~Fulman, Carries, shuffling, and symmetric functions, \emph{Advances in Applied Mathematics} 43 (2009),
176--196.
\item L.~A.~Santaló, \emph{Integral Geometry and Geometric Probability}, Addison-Wesley, 1976.
\item T.~Tao and V.~Vu, On random $\pm1$ matrices: singularity and determinant, \emph{Random Structures \& Algorithms} 28 (2006),
1--23.
\item R.~E.~A.~C.~Paley, On orthogonal matrices, \emph{Journal of Mathematics and Physics} 12 (1933), 311--320; K.~J.~Horadam,
\emph{Hadamard Matrices and Their Applications}, Princeton, 2007.
\item R.~Salem and A.~Zygmund, Some properties of trigonometric series whose terms have random signs, \emph{Acta Mathematica}
91 (1954), 245--301; P.~Balister, B.~Bollobás, R.~Morris, J.~Sahasrabudhe and M.~Tiba, Flat Littlewood polynomials exist,
\emph{Annals of Mathematics} 192 (2020), 977--1004.
\item W.~Rudin, Some theorems on Fourier coefficients, \emph{Proceedings of the American Mathematical Society} 10 (1959),
855--859.
\item H.~Minkowski, Zur Theorie der positiven quadratischen Formen, \emph{Journal für die reine und angewandte Mathematik}
101 (1887), 196--202; M.~Senechal, \emph{Quasicrystals and Geometry}, Cambridge, 1995.
\end{itemize}
